On considère la fonction $f$ définie sur $]0~;~+\infty[$ par :
\[
f(x)=\dfrac{\ln x}{x^2}.
\]

\subsection*{Partie A}

\begin{enumerate}
	\item \begin{enumerate}
			\item Lorsque $x$ tend vers $0$ par valeurs positives, on a :
			\[
			\lim\limits_{x\to 0^+}\ln x=-\infty
			\quad\text{et}\quad
			\lim\limits_{x\to 0^+}x^2=0^+.
			\]
			Comme $x^2>0$ sur $]0~;~+\infty[$, on obtient :
			\[
			\lim\limits_{x\to 0^+} f(x)=-\infty.
			\]
			Lorsque $x$ tend vers $+\infty$, on utilise la croissance comparée :
			\[
			\lim\limits_{x\to +\infty}\dfrac{\ln x}{x^2}=0.
			\]
			Ainsi :
			\[
			\lim\limits_{x\to +\infty} f(x)=0.
			\]

		\item Graphiquement, la droite d'équation $x=0$ est une asymptote verticale à la courbe $\mathcal C_f$.
		De plus, la droite d'équation $y=0$ est une asymptote horizontale à $\mathcal C_f$ au voisinage de $+\infty$.
	\end{enumerate}

	\item Pour tout $x>0$, on peut écrire :
	\[
	f(x)=\ln x\times x^{-2}.
	\]
	Donc :
	\[
	f'(x)=\dfrac1x\times x^{-2}+\ln x\times\left(-2x^{-3}\right).
	\]
	Ainsi :
	\[
	f'(x)=x^{-3}-2\ln x\,x^{-3}=\dfrac{1-2\ln x}{x^3}.
	\]

	\item Comme $x^3>0$ pour tout $x>0$, le signe de $f'(x)$ est celui de $1-2\ln x$.
	On résout :
	\[
	1-2\ln x=0
	\quad\Longleftrightarrow\quad
	\ln x=\dfrac12
	\quad\Longleftrightarrow\quad
	x=\sqrt{\e}.
	\]
	Ainsi $f'(x)>0$ sur $]0~;~\sqrt\e[$ et $f'(x)<0$ sur $]\sqrt\e~;~+\infty[$.
	De plus :
	\[
	f(\sqrt\e)=\dfrac{\ln(\sqrt\e)}{(\sqrt\e)^2}
	=\dfrac{\frac12}{\e}
	=\dfrac{1}{2\e}.
	\]

		On obtient le tableau de variation suivant :

		\begin{center}
		\renewcommand{\arraystretch}{1.4}
		\begin{tabular}{|c|ccc|}
		\hline
		$x$ & $0$ & $\sqrt\e$ & $+\infty$ \\
		\hline
		Signe de $f'(x)$ & $+$ & $0$ & $-$ \\
		\hline
		Variations de $f$ &
		$\lim\limits_{x\to 0^+} f(x)=-\infty$
		& $\dfrac{1}{2\e}$
		& $\lim\limits_{x\to +\infty} f(x)=0$ \\
		\hline
		\end{tabular}
		\end{center}

	\item On cherche l'équation de la tangente à la courbe $\mathcal C_f$ au point d'abscisse 1.
	On a :
	\[
	f(1)=\dfrac{\ln 1}{1^2}=0
	\]
	et
	\[
	f'(1)=\dfrac{1-2\ln 1}{1^3}=1.
	\]
	L'équation de la tangente est donc :
	\[
	y=f'(1)(x-1)+f(1),
	\]
	soit :
	\[
	\Delta : y=x-1.
	\]

	\item Pour tout $x>0$ :
	\[
	f'(x)=(1-2\ln x)x^{-3}.
	\]
	Donc :
	\[
	f''(x)=\left(-\dfrac2x\right)x^{-3}+(1-2\ln x)(-3x^{-4}).
	\]
	Ainsi :
	\[
	f''(x)=-2x^{-4}-3(1-2\ln x)x^{-4}
	=\left(-2-3+6\ln x\right)x^{-4}.
	\]
	Finalement :
	\[
	f''(x)=\dfrac{-5+6\ln x}{x^4}.
	\]

	\item \begin{enumerate}
		\item Comme $x^4>0$ pour tout $x>0$, le signe de $f''(x)$ est celui de $-5+6\ln x$.
		On résout :
		\[
		-5+6\ln x=0
		\quad\Longleftrightarrow\quad
		\ln x=\dfrac56
		\quad\Longleftrightarrow\quad
		x=\e^{\frac56}.
		\]
		Ainsi :
		\[
		f''(x)<0 \text{ sur } ]0~;~\e^{\frac56}[
		\quad\text{et}\quad
		f''(x)>0 \text{ sur } ]\e^{\frac56}~;~+\infty[.
		\]
		La fonction $f$ est donc concave sur $]0~;~\e^{\frac56}]$ et convexe sur $[\e^{\frac56}~;~+\infty[$.

		Le point d'inflexion a pour abscisse $\e^{\frac56}$ et pour ordonnée :
		\[
		f\left(\e^{\frac56}\right)
		=\dfrac{\frac56}{\left(\e^{\frac56}\right)^2}
		=\dfrac{\frac56}{\e^{\frac53}}
		=\dfrac{5}{6\e^{\frac53}}.
		\]
		Le point d'inflexion est donc :
		\[
		\left(\e^{\frac56}~;~\dfrac{5}{6\e^{\frac53}}\right).
		\]

		\item Sur l'intervalle $]0~;~\e^{\frac56}]$, la fonction $f$ est concave.
		Or, pour une fonction concave, la courbe représentative est située en dessous de ses tangentes.
		La tangente à $\mathcal C_f$ au point d'abscisse 1 a pour équation $y=x-1$.
		Donc, pour tout $x\in ]0~;~\e^{\frac56}]$ :
		\[
		f(x)\leqslant x-1.
		\]
		Autrement dit :
		\[
		x-1\geqslant \dfrac{\ln x}{x^2}.
		\]
	\end{enumerate}

	\item Posons $a=\e^{\frac56}$.
	D'après la question précédente, on a :
	\[
	a-1\geqslant f(a).
	\]
	De plus, comme $a=\e^{\frac56}>\sqrt\e$, la fonction $f$ est décroissante sur $[a~;~+\infty[$.
	Ainsi, pour tout $x\geqslant a$ :
	\[
	f(x)\leqslant f(a).
	\]
	Par ailleurs, pour tout $x\geqslant a$ :
	\[
	x-1\geqslant a-1.
	\]
	On en déduit :
	\[
	x-1\geqslant a-1\geqslant f(a)\geqslant f(x).
	\]
	Donc, pour tout $x\in\left[\e^{\frac56}~;~+\infty\right[$ :
	\[
	x-1\geqslant \dfrac{\ln x}{x^2}.
	\]
\end{enumerate}

\subsection*{Partie B}

Pour tout entier naturel non nul $n$, on pose :
\[
I_n=\int_1^n\dfrac{\ln x}{x^2}\,\d x.
\]

\begin{enumerate}
	\item Pour $n\geqslant 1$, l'intégrale $I_n$ représente l'aire, exprimée en unités d'aire, du domaine compris entre la courbe $\mathcal C_f$, l'axe des abscisses, et les droites d'équations $x=1$ et $x=n$.
	En effet, pour $x\geqslant 1$, on a $\ln x\geqslant 0$, donc $f(x)\geqslant 0$.

	\item Pour tout entier naturel non nul $n$ :
	\[
	I_{n+1}-I_n=\int_1^{n+1}f(x)\,\d x-\int_1^n f(x)\,\d x
	=\int_n^{n+1} f(x)\,\d x.
	\]
	Or, pour tout $x\geqslant n\geqslant 1$, on a $f(x)\geqslant 0$.
	Donc :
	\[
	I_{n+1}-I_n\geqslant 0.
	\]
	La suite $(I_n)$ est donc croissante.

	\item Les rectangles représentés ont une largeur égale à 1.
	Le premier rectangle, sur l'intervalle $[1~;~2]$, a pour hauteur le maximum de $f$ sur cet intervalle, c'est-à-dire $f(\sqrt\e)=\dfrac{1}{2\e}$.
	Pour les rectangles suivants, sur les intervalles $[i~;~i+1]$ avec $2\leqslant i\leqslant 9$, la hauteur est $f(i)=\dfrac{\ln i}{i^2}$.

	Un script possible est donc :

	\begin{center}
	\begin{tblr}{colspec={X[8cm,font=\ttfamily]},hline{1,Z}={solid}, vlines}
from math import *\\
\\
S=1/(2*exp(1))\\
for i in range(2,10):\\
\quad S = S + log(i)/(i*i)\\
print(S)
	\end{tblr}
	\end{center}

	\item On calcule $I_n$ par intégration par parties.
	On pose :
	\[
	u(x)=\ln x
	\quad\text{et}\quad
	v'(x)=\dfrac1{x^2}.
	\]
	Alors :
	\[
	u'(x)=\dfrac1x
	\quad\text{et}\quad
	v(x)=-\dfrac1x.
	\]
	Donc :
	\[
	\begin{aligned}
	I_n
	&=\left[-\dfrac{\ln x}{x}\right]_1^n-\int_1^n\left(-\dfrac1x\right)\times\dfrac1x\,\d x\\
	&=-\dfrac{\ln n}{n}+\dfrac{\ln 1}{1}+\int_1^n\dfrac1{x^2}\,\d x\\
	&=-\dfrac{\ln n}{n}+\left[-\dfrac1x\right]_1^n\\
	&=-\dfrac{\ln n}{n}-\dfrac1n+1\\
	&=\dfrac{n-1-\ln n}{n}.
	\end{aligned}
	\]

	\item On a :
	\[
	I_n=\dfrac{n-1-\ln n}{n}=1-\dfrac1n-\dfrac{\ln n}{n}.
	\]
	Or :
	\[
	\lim\limits_{n\to +\infty}\dfrac1n=0
	\quad\text{et}\quad
	\lim\limits_{n\to +\infty}\dfrac{\ln n}{n}=0.
	\]
	Donc :
	\[
	\lim\limits_{n\to +\infty} I_n=1.
	\]
\end{enumerate}
